\chapter{Introdução}

% De acordo com Gil (2018), a introdução se inicia com a apresentação do tema do projeto e do problema que se pretende solucionar com a pesquisa, assim como sua delimitação espacial e temporal. Segundo Brasileiro (2021), a introdução serve para situar o leitor sobre o que está sendo pesquisado, por que, para que e como. 

%O que são os adobes, como técnica construtiva, e por que voltaram a ser discutidos no contexto atual da construção civil? 

Os adobes consistem em uma técnica construtiva de terra crua, em que os blocos são fabricados a partir de uma mistura de argila, areia e água, moldados e secos naturalmente, sem passar por processos de queima, frequentemente recebendo a adição de fibras naturais para aprimorar suas propriedades \cite{Bailly2024, Junior2020, Dormohamadi2024, Morel2021}. Apesar de figurar entre as soluções habitacionais mais antigas da humanidade, o uso do adobe entrou em declínio no pós-Revolução Industrial, sendo substituído por materiais convencionais,como cimento e tijolos cozidos, em virtude de limitações técnicas naturais, sobretudo sua baixa resistência à compressão, baixa durabilidade e alta vulnerabilidade à água e intempéries \cite{Junior2020, Morel2021, Rocco2024, Sharma2016, VanDerLinden2019}.

No entanto, no contexto atual da construção civil, o material tem passado por uma revalorização impulsionada pela necessidade urgente de práticas sustentáveis, redução do impacto ambiental, eficiência de custos e alinhamento à economia circular \cite{Bailly2024, Morel2021, VanDerLinden2019}. Para viabilizar seu uso moderno e corrigir as fraquezas tradicionais, a engenharia tem focado na estabilização física e química dos blocos, utilizando aditivos como o cimento para reduzir a absorção de água e elevar a resistência \cite{Bailly2024, Dormohamadi2024, Sharma2016, Kariyawasam2016}. Contudo, essa intervenção exige controle rigoroso: teores inadequados de cimento podem prejudicar o isolamento térmico e a permeabilidade do bloco, enquanto o tempo de cura se mostra um fator determinante, exigindo extensões de até 90 dias para garantir a secagem completa, evitar a retração por fissuras e consolidar o potencial máximo de resistência à compressão do material \cite{Bailly2024}.

%Quais são as características dos solos da região de Barra do Garças–MT e de que forma essas características condicionam o uso do adobe e a necessidade de estabilização com cimento? 

Os solos da região de Barra do Garças e arredores são caracterizados predominantemente como materiais granulares do tipo areia siltosa ou argilosa (classe AASHTO A-2-4), apresentando cerca de 66,37\% de areia fina e um reduzido Índice de Plasticidade de 6\%  \cite{Barroso2019, Barroso2023, Santos2018, Silva2019}. Essa configuração física, marcada por uma elevada fração arenosa e escassez de finos, inviabiliza a confecção de adobes tradicionais devido à carência de argila para conferir coesão, resultando em blocos que se desintegram facilmente e possuem baixa resistência estrutural \cite{Junior2021, Silva2019}. Por conseguinte, a estabilização com cimento Portland torna-se essencial para viabilizar o uso desse solo, atuando por meio de reações de hidratação que geram uma matriz aglutinante capaz de preencher vazios e ligar as partículas de areia, o que garante o incremento necessário na resistência à compressão e a durabilidade do material perante a exposição hídrica \cite{Novato2019, Silva2019}.

%De que forma o comportamento mecânico dos adobes é influenciado por variáveis de produção e de exposição, como teor de cimento, tempo de cura e presença de água? 

O comportamento mecânico dos adobes estabilizados é intrinsecamente dependente de variáveis de produção, como o teor de cimento, que eleva a rigidez e a resistência à compressão por meio de mecanismos químicos \cite{Roshan2022, Zhang2025, Santos2018, Aqtash2021}, como a formação de géis cimentantes (C-S-H e C-A-H) e reações pozolânicas, e físicos, incluindo a floculação de partículas via trocas catiônicas e o preenchimento de vazios intersticiais \cite{Zhang2025, Ding2021, Messis2022}. Esse ganho de desempenho é potencializado pelo tempo de cura, visto que a transição dos 7 para os 28 dias promove a maturação microestrutural, substituindo agulhas incipientes de hidratação por placas densas que cimentam a matriz do solo e reduzem significativamente a porosidade \cite{Yang2024, Wang2022}. Contudo, a integridade do material é severamente desafiada por variáveis de exposição hídrica; a presença de umidade e a submissão a ciclos de molhagem e secagem desencadeiam processos de expansão e retração dos argilominerais, lixiviação de ligantes e fadiga mecânica. Microscopicamente, essa dinâmica resulta na expansão de poros e na formação de redes de microfissuras que, macroscopicamente, comprometem a estabilidade estrutural e acarretam uma queda abrupta na capacidade de carga e nas resistências ao cisalhamento da alvenaria \cite{Liu2023, Chen2024}.

%De acordo com estudos científicos recentes, como essas variáveis têm sido investigadas em materiais à base de terra? Esses estudos analisam essas variáveis de forma isolada ou integrada? 

Historicamente, a maioria das pesquisas tem avaliado as variáveis de estabilização de forma isolada e aplicada a apenas um ambiente ou tipo de solo específico , mantendo uma forte preferência por análises pontuais de resistência à compressão em detrimento de ensaios de durabilidade frente à exposição severa \cite{Khunaprapakorn2025, Junior2024}. No entanto, analisar esses parâmetros separadamente apresenta limitações técnicas significativas , visto que a evolução da resistência estrutural e os processos químicos de hidratação dependem de uma combinação altamente complexa entre a dosagem de cimento, o tempo de cura e o estado de umidade \cite{Abdallah2023}. Para superar essa barreira, estudos científicos recentes têm migrado para abordagens integradas, combinando investigações macroscópicas e microscópicas para estabelecer correlações precisas e fornecer uma compreensão holística dos mecanismos de estabilização \cite{Yang2024}. Adicionalmente, observa-se o emprego de ferramentas avançadas de análise de dados, como a %\gls{acp}  e \gls{rna}, para o desenvolvimento de modelos que avaliam a interdependência simultânea entre o teor de cimento, o tempo de cura, a densidade e as variações no teor de água \cite{Abdallah2023}.

%Quais limitações ou lacunas são apontadas na literatura quanto à compreensão conjunta desses efeitos?

A literatura técnica estabelece que o aumento do tempo de cura promove o refinamento da microestrutura e a redução da permeabilidade, embora, em baixos teores de cimento, a formação limitada de produtos de hidratação impeça a obstrução completa dos macroporos e das redes de fluxo \cite{Javed2025}. Todavia, a compreensão conjunta desses parâmetros frente à exposição à água permanece um desafio, visto que a maioria das pesquisas analisa as variáveis de forma isolada, falhando em capturar a interdependência entre a dosagem do ligante, a cinética de cura e a resistência residual em serviço \cite{Abdallah2023, Pham2022, Giada2019, Morel2021}. Soma-se a esse cenário a inadequação de ensaios normatizados excessivamente severos e o notável descompasso entre os resultados de laboratório e a performance real em campo, onde as flutuações de umidade e a heterogeneidade do material alteram drasticamente o comportamento mecânico projetado em estado seco \cite{Traore2021, Tarham2024, Giada2019, Khunaprapakorn2025, Morel2021}. Diante da escassez de modelos matemáticos que integrem simultaneamente a história de cura e o dano por imersão, justifica-se a necessidade deste estudo para avaliar como a maturação inicial do solo-cimento mitiga a perda de resistência sob condições de saturação.

%\newpage

\section{Problematização}

% PROBLEMATIZAÇÃO A problematização é a pergunta que irá nortear a pesquisa. O problema deve ser, sobretudo, claro e delimitado para que sua execução se torne viável. 

%De que forma a ação da água compromete o desempenho mecânico e a durabilidade dos adobes? 

A ação da água em materiais terrosos, como o adobe, está associada a mecanismos físicos que controlam sua integridade estrutural e seu desempenho ao longo do tempo. Em materiais não estabilizados, a resistência mecânica está diretamente relacionada à sucção capilar, responsável pela coesão entre partículas. A presença de umidade reduz essa sucção, promovendo o enfraquecimento hidromecânico do material \cite{Gerard2015, Beckett2020, Bui2014}. Além disso, a água é armazenada e transportada no interior do adobe por meio de processos de adsorção em microporos e capilaridade em macroporos, sendo a ascensão capilar e a infiltração por precipitação as principais vias de entrada. A interação contínua com o meio ambiente, associada à natureza higroscópica do material, favorece a ocorrência de ciclos de umedecimento e secagem, que induzem expansão e retração das partículas de argila, gerando tensões internas e fissuração progressiva \cite{Gerard2015, Dieudonne2016, Beckett2020, Bui2014}.

%(GERARD et al., 2015; DIEUDONEE , DELLA VECCHIA, CHARLIER, 2016; SATHIPARAN & RUMESHKUMAR, 2018 BECKETT, JAQUIN & MOREL, 2020; BUI et al., 2021; HART et al., 2021). 

Esses mecanismos físicos desencadeiam processos de degradação que comprometem a durabilidade e o desempenho mecânico dos adobes. A ação da chuva pode provocar erosão superficial por impacto de gotas e escoamento, resultando em perda de material e redução da seção resistente. Simultaneamente, a infiltração prolongada pode levar à saturação do material, reduzindo significativamente sua resistência à compressão e ao cisalhamento, podendo inclusive conduzir à perda de estabilidade estrutural \cite{Gerard2015, Bui2014}. A formação de fissuras decorrente dos ciclos higrotérmicos intensifica esse processo ao criar caminhos preferenciais para a entrada de água. Em condições mais severas, especialmente em materiais não estabilizados, a redução da coesão interna pode levar à desintegração parcial ou total do material, evidenciando que a água atua como o principal agente de degradação em estruturas de adobe \cite{Beckett2020, Ratchakrom2021, Guettatfi2021}. 

% (GERARD et al., 2015; BUI et al., 2021; HART et al., 2021)
%(BECKETT, JAQUIN & MOREL, 2020; GUETTATFI ET AL., 2021;HART et al., 2021; RATCHAKROM & RODVINIJ, 2021).

%Quais limitações estão associadas às estratégias historicamente utilizadas para mitigar esses efeitos? 

Historicamente, a estabilização química com cimento Portland tem sido a principal estratégia para mitigar a baixa resistência e instabilidade volumétrica de solos expansivos, porém essa abordagem apresenta limitações críticas de desempenho e sustentabilidade. Embora o cimento promova um aumento substancial na rigidez mecânica, ele transforma a matriz do solo em um material inerentemente frágil e suscetível à fissuração sob tensão, o que compromete sua integridade estrutural a longo prazo \cite{Chen2024, Zhang2025, Lu2019}. Adicionalmente, a eficácia dessa técnica é severamente prejudicada em solos com presença de sulfatos, onde a formação de minerais expansivos como a etringita pode causar a deterioração química e física da matriz, tornando o uso de dosagens elevadas financeiramente proibitivo e ambientalmente insustentável \cite{Mahedi2018}.

No que tange à durabilidade sob condições climáticas, as estratégias tradicionais falham em garantir estabilidade frente a ciclos de umedecimento e secagem, que destroem progressivamente a coesão do material e podem reduzir a resistência à compressão em mais de 50\% \cite{Chen2024, Hu2024}. Essa vulnerabilidade é agravada por limitações nos processos de cura; a exposição precoce à água ou a altas temperaturas pode induzir o "efeito de cruzamento" (\textit{cross-over effect}), resultando em uma microestrutura de poros grosseiros e resistência final a longo prazo inferior à esperada \cite{AlGburi2025, Wang2022, Lu2019}. Portanto, a dependência de altos teores de ligantes para conter a retração e o intemperismo evidencia a necessidade de compreender as correlações entre o tempo de cura e a exposição hídrica para otimizar o desempenho desses materiais estabilizados \cite{Zhang2025}.

%Como essas limitações, associadas ao processo de industrialização da construção civil, contribuíram para a redução do uso do adobe ao longo do tempo? 

O declínio do uso do adobe na construção civil moderna está intrinsecamente ligado às suas limitações mecânicas e à alta vulnerabilidade hídrica, que o colocam em desvantagem perante os materiais industrializados. O material apresenta baixa resistência à tração e comportamento frágil \cite{Brito2023, Fages2022}, além de uma resistência à compressão que, embora aceitável no estado seco (variando entre 0,36 MPa e 2,15 MPa), reduz-se drasticamente em condições de saturação \cite{Kafodya2019, Khunaprapakorn2025, Dao2018}. Essa extrema sensibilidade à umidade resulta em patologias graves, como a erosão severa, frequentemente descrita como o "derretimento" da estrutura, e instabilidades estruturais causadas por ciclos de umedecimento e secagem \cite{Day1993, Sharma2016}. Tais fragilidades técnicas são agravadas pela ausência de normas técnicas e códigos de prática modernos, o que gera insegurança jurídica e técnica para projetistas, dificultando a aprovação de projetos e a garantia de segurança das edificações \cite{Fages2022, Martins2018, Morel2021, Victor2025}.

Paralelamente aos obstáculos técnicos, o processo de industrialização consolidou uma "ideologia do progresso" que marginalizou o adobe ao associá-lo a um status inferior e à precariedade \cite{Aspillaga2024, Morel2021, Novato2019}. A transição para materiais "duros", como o concreto e o aço, transformou a construção com terra em um símbolo de privação material, rotulando-a erroneamente como o "bloco de cimento dos pobres" \cite{Zoungrana2021}. Esse estigma social, alimentado pela falta de conhecimento técnico e por preocupações históricas com a higiene, fez com que o adobe fosse percebido como uma solução arcaica, adequada apenas a contextos rurais ou populações de baixa renda \cite{Soler2021, Novato2019, Morel2021, Aspillaga2024}. Consequentemente, a busca pela "modernidade" urbana e pelo prestígio social atrelado aos materiais industrializados acelerou a substituição das técnicas tradicionais, restringindo o uso do adobe e dificultando sua integração em práticas contemporâneas de construção sustentável \cite{Zoungrana2021}.

%Embora existam avanços recentes (impermeabilizantes, hidrofugantes, aditivos, etc.), em que medida essas soluções são suficientes e quais incertezas ainda permanecem? 

A utilização de avanços tecnológicos, como emulsões hidrofóbicas e agentes superhidrofóbicos, tem demonstrado eficácia na preservação da integridade estrutural de solos estabilizados sob condições de saturação prolongada, permitindo que matrizes tratadas mantenham sua resistência por até 30 dias de imersão, enquanto amostras convencionais colapsam em poucas horas \cite{Guo2025}. Contudo, a suficiência dessas soluções é questionada devido ao impacto negativo na resistência à compressão simples, especialmente em matrizes de cimento Portland. Essa incerteza mecânica decorre da inibição das reações do silicato tricálcico ($C_3S$) e do retardamento da hidratação, que resultam em uma microestrutura porosa e um "esqueleto" sólido frouxo, o que pode comprometer o desempenho estrutural nas idades iniciais \cite{Lu2023, Luo2023, Jahandari2023}.

Adicionalmente, embora a literatura sugira que aditivos densificadores, nanomateriais e materiais pozolânicos possam atuar como substitutos parciais do cimento para reduzir custos e emissões de $CO_2$, a estabilização de longo prazo permanece dependente de uma dosagem equilibrada \cite{Zhang2025, Jahandari2023, Dormohamadi2024, Wahabi2021}. A redução drástica do teor de cimento, sem o suporte de uma rede de hidratação consolidada ou reforço de fibras, pode limitar a eficácia da estabilização apenas ao curto prazo, tornando o material suscetível à desintegração precoce sob ciclos de umedecimento e secagem \cite{Wahabi2021}. Portanto, a lacuna de conhecimento reside em determinar até que ponto o aumento do teor de cimento pode ser tecnicamente substituído por aditivos químicos sem que a durabilidade e a capacidade de suporte sejam sacrificadas em ambientes de alta exposição à água.

%Como a interação entre teor de cimento, tempo de cura e condição de exposição à água afeta a resistência à compressão dos adobes, e qual o papel dos hidrofugantes nesse contexto?

A estabilização de adobes via cimentação depende intrinsecamente da evolução das reações de hidratação e da formação de silicatos de cálcio hidratado (C-S-H), que consolidam a matriz e reduzem a permeabilidade hídrica \cite{Dao2018}. No entanto, a eficácia dessa barreira é sensível ao fator tempo; curas precoces ou insuficientes (inferiores a 3 dias) podem levar a falhas estruturais sob ciclos de umedecimento, mesmo em traços com elevado teor de cimento \cite{Wahabi2021, Nshimiyimana2020, Yang2024}. Embora a literatura aponte para um limiar técnico-econômico ótimo, frequentemente associado a 6\% de cimento e 7 dias de cura para solos específicos, a ausência de um consenso universal, decorrente da variabilidade mineralógica e granular dos solos, impõe desafios à padronização do desempenho mecânico sob imersão \cite{Abdallah2023, Wahabi2021}.

Nesse contexto, o uso de hidrofugantes surge como uma alternativa para mitigar a tensão capilar e a retração autógena, sem a necessidade de onerar o traço com ligantes hidráulicos adicionais. Contudo, sua aplicação não é isenta de riscos, podendo inibir a hidratação inicial e fragilizar a microestrutura nos primeiros sete dias \cite{Lu2023}. A complexidade dessas interações físico-químicas, aliada à carência de dados sobre a durabilidade de longo prazo em condições reais de exposição e modelos que integrem a energia de compactação à evolução hidromecânica, define a lacuna científica que justifica este estudo \cite{Abdallah2023, Jahandari2023, Rosicki2022}. Portanto, torna-se imperativo investigar como a manipulação conjunta do teor de cimento, tempo de cura e aditivos repelentes de água pode garantir a resistência residual necessária para a viabilidade do adobe como material estrutural.

\section{Justificativa}

% JUSTIFICATIVA A justificativa é um texto dissertativo-argumentativo em que o pesquisador deve convencer a comunidade acadêmico científica da relevância da sua proposta. O texto pode apresentar os fatores que influenciaram a escolha do tema e sua relação com a experiência acadêmica ou profissional do autor (BRASILEIRO, 2021; GIL, 2018). 

%Quais propriedades dos adobes (ambientais, térmicas e mecânicas) justificam seu interesse na construção civil, e em que medida sua aplicação depende do desempenho mecânico e da durabilidade do material? 

O adobe destaca-se na construção civil contemporânea por sua baixa pegada de carbono e reduzida energia incorporada, promovendo a sustentabilidade através do uso de solo local e da valorização de resíduos sólidos e agroindustriais \cite{Arduin2022, Khunaprapakorn2025, Arduin2024, Aspillaga2024}. Aliado ao baixo impacto ambiental, sua elevada inércia térmica garante o conforto passivo e a eficiência energética das edificações, agindo na regulação das variações de temperatura interna através de atrasos térmicos e trocas de calor latente por evaporação e condensação \cite{Christoforou2016, Heracleous2023, Rincon2019, Khunaprapakorn2025, Giada2019}. Contudo, a viabilidade técnica do material é severamente condicionada à sua fragilidade mecânica intrínseca e à extrema vulnerabilidade hídrica, uma vez que a exposição contínua à água acarreta rápida erosão e perda drástica da resistência à compressão \cite{Brito2023, Dormohamadi2020, Fages2022,Dao2018,Day1993, Dormohamadi2024}. Por essa razão, a aplicação segura e duradoura do adobe na engenharia moderna depende fundamentalmente do aprimoramento de suas propriedades via técnicas de estabilização, como a adição de cimento e fibras, visando atender aos requisitos estruturais normativos e garantir a resistência à degradação hídrica.

%Qual é a influência do teor de cimento e do tempo de cura no desenvolvimento da resistência mecânica de materiais solo-cimento? 

A utilização do solo-cimento justifica-se por suas expressivas vantagens ambientais e econômicas, destacando-se a ausência do processo de queima e uma demanda energética até 11 vezes menor que a da alvenaria cerâmica, resultando em emissões de apenas 0,07 kg de CO2-eq/kg \cite{Barroso2019, Brito2023, Ibrahim2020}. Sob a ótica socioeconômica, a técnica possibilita reduções de custos entre 30\% e 50\% no valor total da obra, consolidando-se como uma alternativa viável para habitações de interesse social ao superar os requisitos de resistência estipulados por normas técnicas, como a NBR 16814, frequentemente atingindo valores acima de 3,5 MPa \cite{Ibrahim2020, Kariyawasam2016, Junior2020, Dao2018, Bogas2019}. Contudo, a eficiência mecânica da estabilização depende de variáveis intrínsecas ao material, como a porosidade e a mineralogia do solo, visto que dosagens excessivas de aglomerante podem resultar em perda de ductilidade e falhas frágeis indesejadas. Nesse sentido, a determinação do %\gls{teo} por meio do controle do teor de cimento e do tempo de cura é fundamental para garantir o máximo desempenho estrutural com o menor investimento financeiro e impacto ambiental possível \cite{Santana2021}.

%Como a presença de água afeta os mecanismos de degradação e a resistência mecânica de materiais à base de terra, e quais evidências existem sobre a eficácia de tratamentos hidrofugantes na mitigação desses efeitos? 

A presença de água é o principal agente degradador de materiais à base de terra, atuando por mecanismos microestruturais como o \textit{slaking}, que rompe agregados via gradientes de pressão, e ciclos de inchamento e retração que geram tensões internas de sucção. Tais processos promovem a dissolução de ligantes e a coalescência de microfissuras, resultando em quedas de resistência mecânica de até 45\%, impacto que é significativamente acentuado por períodos insuficientes de cura inicial \cite{Hu2018, Liu2023, Ralf2025}. Para mitigar esses efeitos, tratamentos de superfície com silano-siloxano mostram-se eficazes em reduzir a absorção capilar para cerca de 15\% em relação ao controle, embora atuem apenas no bloqueio externo sem refinar a microestrutura interna \cite{Bogas2020}. Por outro lado, o uso de hidrofugantes integrais, como o Caltite, modifica a estrutura física da matriz através de glóbulos poliméricos que bloqueiam vazios capilares e preservam a integridade mecânica sob imersão \cite{Ezreig2022}. Entretanto, a durabilidade dessas proteções é vulnerável em ambientes de alto pH, onde a hidrólise de ligações químicas por íons hidroxila pode converter a superfície repelente em hidrofílica, comprometendo a resistência à compressão a longo prazo \cite{Guo2025}.

%De que forma investigações experimentais controladas podem contribuir para o avanço do conhecimento e para a aplicação mais segura e tecnicamente fundamentada de materiais à base de terra na construção civil?

A viabilidade técnica e a segurança estrutural de materiais à base de terra permanecem condicionadas à superação da heterogeneidade inerente às matérias-primas e à carência de uma base normativa universal, fatores que geram incertezas críticas quanto à durabilidade frente a oscilações de umidade em serviço \cite{Giada2019, Pham2022, Santana2021, Martins2018, Bogas2019, VanDerLinden2019}. Nesse contexto, investigações experimentais controladas mostram-se indispensáveis, pois evidenciam que o ajuste preciso do teor de cimento é o que define a integridade da matriz; enquanto dosagens insuficientes resultam em degradação severa e perda de capacidade de carga, teores ótimos (entre 6\% e 10\%) promovem a blindagem dos poros por meio da formação de produtos de hidratação e reações pozolânicas, garantindo a manutenção da resistência mesmo sob ciclos severos de exposição à água \cite{Morel2021, Martins2018, Roshan2022, Ralf2025, Wahabi2021}. Complementarmente, o controle rigoroso do tempo de cura atua como fator determinante para a maturidade dessas reações, evitando a lixiviação precoce do cálcio e promovendo o adensamento microestrutural necessário para a estabilidade de longo prazo \cite{Ezreig2022,Morel2021}. Portanto, este trabalho justifica-se ao converter lacunas de conhecimento em parâmetros técnicos mensuráveis, estabelecendo a conexão necessária entre o comportamento laboratorial e o desempenho em campo, única via para uma aplicação fundamentada e segura na construção civil contemporânea.

%\section{Objetivos}

% OBJETIVOS A apresentação dos objetivos da pesquisa deve ser feita em termos claros e precisos. Recomenda-se, portanto, que em sua redação sejam utilizados verbos de ação, como avaliar, analisar, descrever, identificar e verificar. 

%\subsection{Objetivo geral}

% Objetivo geral O objetivo geral define o que o pesquisador pretende atingir com sua investigação 

%\subsection{Objetivos específicos}

% Objetivos específicos 285 Os objetivos específicos definem as etapas do trabalho a serem realizadas para que se alcance o objetivo geral.